\section{Υπολογιστικό νέφος και μοντέλα υπηρεσιών}

Το Κεφάλαιο~1 οριοθέτησε το πρόβλημα, τους στόχους και τα ερευνητικά
ερωτήματα της εργασίας. Το παρόν κεφάλαιο δεν επαναλαμβάνει εκείνη την
εισαγωγή, αλλά θεμελιώνει τις έννοιες που απαιτούνται για να αιτιολογηθούν
οι τεχνικές επιλογές των επόμενων κεφαλαίων: τα μοντέλα παροχής υπηρεσιών,
οι ιδιότητες μιας αξιόπιστης υποδομής, οι σχετικές υπηρεσίες AWS και η θέση
της εφαρμογής Spring Boot μέσα σε αυτό το περιβάλλον.

Η ανάγκη φιλοξενίας μιας διαδικτυακής εφαρμογής δεν περιορίζεται στην εύρεση
ενός υπολογιστή που θα παραμένει ενεργός. Ο υπολογιστής πρέπει να είναι
προσβάσιμος με σταθερό τρόπο, να προστατεύεται από μη αναγκαία εισερχόμενη
κίνηση, να διατηρεί τα δεδομένα του μετά από επανεκκίνηση και να μπορεί να
ενημερώνεται χωρίς να διακόπτεται αδικαιολόγητα η λειτουργία του. Το
υπολογιστικό νέφος συγκεντρώνει αυτές τις δυνατότητες σε ένα περιβάλλον όπου
οι πόροι παρέχονται κατ' απαίτηση, προσπελαύνονται μέσω δικτύου και μπορούν να
παρακολουθούνται και να χρεώνονται ανάλογα με τη χρήση. Τα χαρακτηριστικά
αυτά αποτελούν μέρος του ορισμού του NIST για το cloud computing
\cite{mell2011nist}.

Στην πράξη, ο όρος \emph{cloud} δεν περιγράφει μία μοναδική τεχνική λύση.
Περιλαμβάνει διαφορετικούς βαθμούς ευθύνης για τον πάροχο και τον χρήστη. Στο
μοντέλο Infrastructure as a Service (IaaS), ο πάροχος διαθέτει το φυσικό
κέντρο δεδομένων, τη δικτύωση, τον hypervisor και τους εικονικούς πόρους. Ο
χρήστης εξακολουθεί να επιλέγει λειτουργικό σύστημα, λογισμικό εκτέλεσης,
κανόνες πρόσβασης και τρόπο αποθήκευσης των δεδομένων. Στο Platform as a
Service (PaaS) ο πάροχος αναλαμβάνει μεγαλύτερο μέρος του runtime και της
λειτουργικής συντήρησης, ενώ στο Software as a Service (SaaS) ο τελικός
χρήστης καταναλώνει έτοιμη εφαρμογή. Η παρούσα εργασία κινείται κυρίως στο
IaaS: η Amazon EC2 παρέχει την εικονική μηχανή, αλλά η εγκατάσταση Docker, η
MySQL, το Spring Boot, ο Nginx και η διαδικασία ενημέρωσης παραμένουν ευθύνη
του διαχειριστή.

Η διάκριση αυτή είναι σημαντική επειδή αποτρέπει μία συνηθισμένη
παρερμηνεία: η μεταφορά μιας εφαρμογής σε cloud υποδομή δεν την καθιστά
αυτόματα ανεκτική σε βλάβες ή ελαστική. Η EC2 instance της εργασίας είναι
ένας απομακρυσμένος και ελεγχόμενος υπολογιστικός κόμβος. Η σταθερή δημόσια
διεύθυνση, το HTTPS και η αυτοματοποιημένη εκκίνηση των containers βελτιώνουν
τη λειτουργία του, χωρίς όμως να εξαλείφουν το μοναδικό σημείο αστοχίας που
δημιουργεί ο ένας κόμβος. Η διαπίστωση συμφωνεί με τη βιβλιογραφία της
μετανάστευσης στο νέφος, η οποία αντιμετωπίζει τη μετάβαση ως σύνολο
αρχιτεκτονικών αποφάσεων και όχι ως απλή αντιγραφή λογισμικού
\cite{jamshidi2013migration}.

\section{Αρχές σχεδιασμού αξιόπιστων υποδομών}

Για ένα πληροφοριακό σύστημα κτηνιατρικής κλινικής, η αξιοπιστία αποκτά
συγκεκριμένο περιεχόμενο. Ο χρήστης πρέπει να μπορεί να βρει το ραντεβού, τον
ιδιοκτήτη και το ιστορικό του ζώου όταν τα χρειάζεται, ενώ μία ενημέρωση της
εφαρμογής δεν πρέπει να διαγράφει τη βάση ή τα μεταφορτωμένα αρχεία. Στη
θεωρία των dependable systems, η διαθεσιμότητα, η αξιοπιστία, η ακεραιότητα
και η δυνατότητα συντήρησης είναι διακριτές αλλά αλληλένδετες ιδιότητες
\cite{avizienis2004dependability}. Για τον λόγο αυτό, η εργασία δεν χρησιμοποιεί
τον όρο ``αξιόπιστο'' μόνο επειδή μία σελίδα άνοιξε επιτυχώς μία φορά.

Πρώτη αρχή είναι ο διαχωρισμός του προσωρινού runtime από τη μόνιμη
κατάσταση. Ένα container μπορεί να διαγραφεί και να δημιουργηθεί εκ νέου,
ενώ τα δεδομένα της MySQL πρέπει να παραμείνουν σε persistent volume. Το ίδιο
ισχύει για τα ιατρικά αρχεία που μεταφορτώνουν οι χρήστες. Δεύτερη αρχή είναι
η ελαχιστοποίηση της δημόσιας επιφάνειας: προς το Διαδίκτυο είναι αναγκαίο να
εκτίθενται οι θύρες HTTP/HTTPS του Nginx, όχι όμως η MySQL ή η εσωτερική θύρα
του Spring Boot. Τρίτη αρχή είναι η επαληθεύσιμη ενημέρωση. Κάθε νέο πακέτο
πρέπει να ελέγχεται, να μεταφέρεται με ασφαλή τρόπο, να ενεργοποιείται με
περιορισμένη αλλαγή και να ακολουθείται από ελέγχους υγείας.

Η αξιοπιστία συνδέεται επίσης με την παρατηρησιμότητα. Ένας κωδικός HTTP 200
επιβεβαιώνει ότι η δημόσια διαδρομή λειτουργεί τη στιγμή του ελέγχου, αλλά
δεν εξηγεί μόνος του αν η βάση είναι υγιής ή αν εξαντλείται η μνήμη. Για αυτό
χρειάζονται διαφορετικά επίπεδα τεκμηρίων: EC2 status checks, κατάσταση
containers, health check της MySQL, εφαρμοστικά logs και μετρικές CPU ή
δικτύου. Το AWS Well-Architected Framework οργανώνει αντίστοιχες ερωτήσεις
σε έξι πυλώνες και χρησιμοποιείται στην εργασία ως πλαίσιο κριτικής
αξιολόγησης, όχι ως πιστοποίηση της εγκατάστασης
\cite{aws2024wellarchitected}.

\section{Υπηρεσίες Amazon Web Services}

Οι υπηρεσίες AWS που εμφανίζονται στην υλοποίηση δεν επιλέχθηκαν ως
ανεξάρτητη συλλογή προϊόντων. Κάθε μία καλύπτει μία πρακτική ανάγκη της ίδιας
διαδρομής: εκτέλεση, μόνιμη αποθήκευση, ελεγχόμενη πρόσβαση, σταθερή
διευθυνσιοδότηση, επίλυση ονόματος και παρακολούθηση. Η μεταξύ τους σχέση
είναι σημαντικότερη από τη μεμονωμένη περιγραφή τους.

\subsection{Υπολογιστικοί πόροι και κλιμάκωση}

Η Amazon EC2 παρέχει εικονικές μηχανές, οι οποίες στην ορολογία της AWS
ονομάζονται instances \cite{aws2026ec2}. Ο τύπος instance ορίζει έναν συνδυασμό εικονικών
επεξεργαστών, μνήμης και δικτυακών δυνατοτήτων. Η επιλογή
\texttt{t3.small} της εργασίας δεν αποτελεί ισχυρισμό ότι αυτή είναι η
καλύτερη διαμόρφωση για κάθε κλινική. Είναι μία συνειδητή ισορροπία για το
συγκεκριμένο πείραμα, όπου στον ίδιο κόμβο εκτελούνται JVM, MySQL και Nginx.
Η καταλληλότητά της κρίνεται αργότερα με μετρήσεις και όχι μόνο από τα
ονομαστικά χαρακτηριστικά της.

Η οριζόντια κλιμάκωση θα απαιτούσε περισσότερα instances και έναν load
balancer, καθώς και απομάκρυνση της βάσης και των αρχείων από τον τοπικό
κόμβο. Αντίστοιχα, μία managed βάση όπως το Amazon RDS θα μείωνε μέρος της
διαχειριστικής εργασίας. Αυτές οι λύσεις παρουσιάζονται ως εναλλακτικές και
όχι ως υλοποιημένες δυνατότητες, επειδή η βασική εγκατάσταση έχει διαφορετικό
στόχο: να είναι οικονομική, κατανοητή και πλήρως ελέγξιμη.

\subsection{Δικτύωση και επίλυση ονομάτων}

Το Virtual Private Cloud (VPC) αποτελεί το λογικό δίκτυο μέσα στο οποίο
τοποθετείται η EC2 instance. Το public subnet διαθέτει διαδρομή προς Internet
Gateway, ενώ το security group λειτουργεί ως stateful εικονικό firewall στο
επίπεδο του instance \cite{aws2026vpc,aws2026securitygroups}. Στη συγκεκριμένη εγκατάσταση επιτρέπονται δημόσια οι
θύρες 80 και 443. Η θύρα 22 για SSH επιτρέπεται μόνο από μία διεύθυνση
\texttt{/32}. Επομένως, μία αλλαγή της δημόσιας IP του διαχειριστή επηρεάζει
την απομακρυσμένη συντήρηση, όχι την πρόσβαση των επισκεπτών στον ιστότοπο.

Η Elastic IP δίνει σταθερή δημόσια IPv4 διεύθυνση στην instance. Πάνω σε
αυτή βασίζονται οι εγγραφές A του Route~53 για το κύριο domain και το
\texttt{www} \cite{aws2026route53}. Το DNS μετατρέπει το κατανοητό από τον άνθρωπο όνομα σε
διεύθυνση δικτύου, ενώ το TLS certificate επιτρέπει στον browser να
επαληθεύσει το όνομα και να δημιουργήσει κρυπτογραφημένη σύνδεση. Πρόκειται
για διαδοχικά βήματα της ίδιας πρόσβασης και όχι για ανταλλάξιμες υπηρεσίες.

\subsection{Σχεσιακές βάσεις δεδομένων και αποθήκευση}

Το Elastic Block Store (EBS) παρέχει block storage που προσαρτάται στην EC2
instance \cite{aws2026ebs}. Το root volume τύπου \texttt{gp3} φιλοξενεί το λειτουργικό σύστημα,
τα Docker volumes και τα αρχεία της εφαρμογής. Τα δεδομένα της MySQL
διατηρούνται σε named volume, ώστε η επαναδημιουργία του container να μη
συνεπάγεται επαναδημιουργία της βάσης. Η επιλογή αυτή προστατεύει από τον
συνηθισμένο κύκλο ενημέρωσης, όχι από κάθε πιθανή απώλεια. Ένα σφάλμα στο
volume ή μία λανθασμένη διαγραφή εξακολουθούν να απαιτούν ανεξάρτητο
αντίγραφο ασφαλείας.

Η διάκριση μεταξύ volume και snapshot είναι επομένως ουσιαστική. Το volume
είναι ο ενεργός χώρος εργασίας, ενώ το snapshot είναι αντίγραφο σε
συγκεκριμένη χρονική στιγμή. Η στρατηγική αντιγράφων ασφαλείας αξιολογείται
χωριστά στο Κεφάλαιο~4, ώστε η απλή χρήση persistent storage να μην
παρουσιαστεί ως πλήρης πολιτική προστασίας δεδομένων.

\subsection{Παρακολούθηση και παρατηρησιμότητα}

Το Amazon CloudWatch συλλέγει βασικές μετρικές της EC2, όπως χρήση CPU και
κίνηση δικτύου, και μπορεί να υποστηρίξει alarms \cite{aws2026cloudwatch}. Οι μετρικές υποδομής πρέπει
να συνδυάζονται με εφαρμοστικούς ελέγχους, επειδή μία instance μπορεί να
λειτουργεί ενώ μία εσωτερική υπηρεσία έχει αποτύχει. Στην εργασία τα
CloudWatch γραφήματα και οι ειδοποιήσεις κόστους θα παρουσιαστούν μόνο μετά
την τελική τους ρύθμιση. Η διάκριση ανάμεσα σε διαθέσιμη δυνατότητα και σε
επαληθευμένη λειτουργία διατηρείται σε όλα τα επόμενα κεφάλαια.

\section{Spring Boot και εφαρμογές νέφους}

Το Spring Boot χρησιμοποιείται για τη δημιουργία του backend \cite{spring2026boot}. Παρέχει
ενσωματωμένο web runtime, μηχανισμό παραμετροποίησης και ενοποίηση με Spring
Web, Spring Security, Jakarta Persistence και validation. Στο Happy Tails το
backend εκθέτει REST endpoints που ανταλλάσσουν JSON με τη διεπαφή React
\cite{fielding2000rest,react2026docs}. Οι
controllers παραλαμβάνουν το HTTP αίτημα, οι application services εφαρμόζουν
τη λειτουργική λογική και τα repositories αναλαμβάνουν την πρόσβαση στη
MySQL. Ο διαχωρισμός αυτός βοηθά να μη μεταφέρεται η επιχειρησιακή λογική
ούτε στον browser ούτε απευθείας στο επίπεδο αποθήκευσης.

Η εφαρμογή είναι οργανωμένη σε δύο Maven modules. Το
\texttt{pet-clinic-data} περιλαμβάνει entities, repositories και υπηρεσίες
δεδομένων, ενώ το \texttt{pet-clinic-web} περιλαμβάνει το web/API επίπεδο,
την ασφάλεια και το frontend build. Η παραγωγή του τελικού JAR ενσωματώνει τα
στατικά αρχεία που δημιουργεί το Vite. Έτσι, ο Nginx χρειάζεται να προωθεί
την κίνηση σε μία εσωτερική Spring υπηρεσία, ενώ η React εφαρμογή και το API
παρουσιάζονται κάτω από το ίδιο HTTPS origin.

Το frontend είναι single-page application σε React και TypeScript
\cite{react2026docs}. Το
React Router αντιστοιχίζει τις διαδρομές σε σελίδες, το TanStack Query
διαχειρίζεται server state και invalidation μετά από μεταβολές, ενώ το
\texttt{i18next} παρέχει ελληνικό και αγγλικό περιεχόμενο. Το FullCalendar
χρησιμοποιείται για την εβδομαδιαία απεικόνιση ραντεβού και το Recharts για
τη γραφική παράσταση του βάρους. Τα εργαλεία αυτά εξυπηρετούν διαφορετικές
ευθύνες· δεν αντικαθιστούν το backend, το οποίο παραμένει η έγκυρη πηγή των
δεδομένων.

Η εκτέλεση σε container βοηθά να μεταφερθεί η ίδια έκδοση του runtime από το
περιβάλλον build στο EC2 instance. Δεν καταργεί όμως την ανάγκη
παραμετροποίησης. Τα database credentials, το JWT signing secret, τα OAuth
credentials και τα URLs αλλάζουν ανά περιβάλλον και παρέχονται μέσω
environment variables. Αυτή η διάκριση κώδικα και ρύθμισης αποτελεί βασική
προϋπόθεση για επαναλήψιμο deployment.

\section{Σχετικές προσεγγίσεις και θέση της εργασίας}

Οι περισσότερες γενικές αναφορές στο cloud δίνουν έμφαση στην ελαστικότητα,
στην οικονομία κλίμακας και στην ταχεία παροχή πόρων
\cite{armbrust2010view}. Οι αρχές αυτές είναι χρήσιμες, αλλά ένα μικρό
πληροφοριακό σύστημα έχει συχνά διαφορετική αφετηρία: περιορισμένο φορτίο,
μικρή ομάδα διαχείρισης και ανάγκη να είναι κατανοητό ολόκληρο το μονοπάτι
από τον browser έως τη βάση. Για αυτό η εργασία εξετάζει πρώτα μία
containerized εγκατάσταση ενός κόμβου και στη συνέχεια συζητά τι θα άλλαζε
με managed database, load balancer ή περισσότερες ζώνες διαθεσιμότητας.

Η εφαρμογή λειτουργεί ταυτόχρονα ως πληροφοριακό σύστημα και ως workload της
αξιολόγησης. Δεν χρησιμοποιείται ένα τεχνητό endpoint που επιστρέφει σταθερή
απάντηση. Οι έλεγχοι ακολουθούν ρεαλιστικές ροές: σύνδεση, καταχώριση
ιδιοκτήτη και ζώου, ραντεβού, ενημέρωση ιατρικού ιστορικού και ανάκτηση των
δεδομένων. Με αυτόν τον τρόπο οι ιδιότητες της υποδομής συσχετίζονται με
λειτουργίες που έχουν νόημα για τον τελικό χρήστη.

\section{Σύνοψη κεφαλαίου}

Το κεφάλαιο τοποθέτησε τις τεχνολογίες της εργασίας στο κατάλληλο πλαίσιο.
Η EC2 καλύπτει την εκτέλεση, το EBS τη μόνιμη block αποθήκευση, το security
group τον έλεγχο εισερχόμενης κίνησης, η Elastic IP τη σταθερή δημόσια
διεύθυνση και το Route~53 την επίλυση του domain. Πάνω σε αυτή τη βάση, ο
Nginx, το Spring Boot, η MySQL και το React σχηματίζουν την εφαρμοστική
στοίβα. Το επόμενο κεφάλαιο εξετάζει πώς αυτή η στοίβα εξυπηρετεί τις
συγκεκριμένες ανάγκες του Happy Tails πριν παρουσιαστούν οι λεπτομέρειες της
cloud υποδομής.
